% ============================================
% MODULAÇÃO ANGULAR: FM E PM
% ============================================

\section{Introdução}

\begin{frame}{Modulação angular versus AM}

\small\compacto
\parbox[t][1.1cm]{\linewidth}{%
\only<1>{Mensagem $m(t)$: um tom lento. Onde vamos escrevê-lo na portadora?}%
\only<2>{AM: a mensagem vai para a amplitude, e a oscilação mantém o ritmo.
  \[ s_{\AM}(t) = A_c[1+\mu\, m(t)]\cos(2\pi f_c t) \]}%
\only<3>{Modulação angular: a amplitude fica fixa, e a mensagem acelera ou freia a oscilação.
  \[ s(t) = A_c\cos[\theta(t)] = A_c\cos[2\pi f_c t + \phi(t)] \]}%
\only<4>{Mesma amplitude o tempo todo: a informação está só nos cruzamentos por zero.
  \[ s(t) = A_c\cos[2\pi f_c t + \phi(t)] \]}}

\begin{center}
\begin{tikzpicture}
\begin{axis}[onda, height=1.1cm, domain=0:2, ymin=-1.3, ymax=1.3, ylabel={$m(t)$}]
\addplot[thick, cBlue] {cos(deg(2*pi*x))};
\end{axis}
\onslide<2->{
\begin{axis}[onda, height=1.65cm, domain=0:2, ymin=-1.8, ymax=1.8, at={(0,-2.3cm)}, ylabel={$s(t)$}]
\only<2>{
  \addplot[cRed, samples=700] {(1+0.6*cos(deg(2*pi*x)))*cos(deg(2*pi*10*x))};
  \addplot[cRed, dashed] {1+0.6*cos(deg(2*pi*x))};
  \addplot[cRed, dashed] {-(1+0.6*cos(deg(2*pi*x)))};
}
\only<3->{
  \addplot[cGreen, samples=900] {cos(deg(2*pi*8*x + 5*sin(deg(2*pi*x))))};
  \addplot[gray, dashed, domain=0:2] {1.15};
  \addplot[gray, dashed, domain=0:2] {-1.15};
}
\only<4>{
  \draw[cRed, dashed] (axis cs:1,-1.8) -- (axis cs:1,1.8);
  \node[nota=cRed, anchor=south] at (axis cs:1,1.25) {$m$ máximo: oscilação mais rápida};
  \draw[cBlue, dashed] (axis cs:0.5,-1.8) -- (axis cs:0.5,1.8);
  \draw[cBlue, dashed] (axis cs:1.5,-1.8) -- (axis cs:1.5,1.8);
  \node[nota=cBlue, anchor=south west] at (axis cs:1.52,1.25) {$m$ mínimo: oscilação mais lenta};
}
\end{axis}}
\end{tikzpicture}
\end{center}

\vspace{-0.45cm}
\onslide<4>{Duas variantes: em FM, a frequência instantânea segue $m(t)$; em PM, a fase segue $m(t)$.}

\end{frame}

% ============================================

\begin{frame}{Fase e frequência instantâneas}

\small
\begin{columns}[T]
\column{0.44\textwidth}
Para $s(t) = A_c \cos[\theta(t)]$:

\vspace{0.15cm}
Fase instantânea
\[
\theta(t) = 2\pi f_c t + \phi(t)
\]

\onslide<3->{Frequência instantânea
\[
f_i(t) = \tikzmarknode{fi2pi}{\frac{1}{2\pi}}\,\tikzmarknode{fidtheta}{\frac{d\theta(t)}{dt}}
\]
\vspace{0.1cm}
\[
f_i(t) = f_c + \frac{1}{2\pi}\frac{d\phi(t)}{dt}
\]}

\column{0.56\textwidth}
\begin{tikzpicture}
\begin{axis}[width=7.6cm, height=5.4cm, axis lines=left, domain=0:4, samples=200,
             xmin=0, xmax=4.2, ymin=-1, ymax=29, xtick=\empty, ytick=\empty,
             xlabel={$t$}, ylabel={$\theta(t)$}, clip=false,
             every axis x label/.style={at={(current axis.right of origin)}, anchor=west, font=\scriptsize},
             every axis y label/.style={at={(current axis.north west)}, anchor=south, font=\scriptsize}]
\onslide<2->{
  \addplot[gray, dashed] {2*pi*x} node[pos=0.97, above left, font=\scriptsize, gray] {$2\pi f_c t$};
  \addplot[very thick, cBlue] {2*pi*x + 1.8*sin(deg(pi*x))};
}
\only<4->{
  \addplot[cRed, thick, domain=1.75:2.25] {4*pi + 3.8*pi*(x-2)};
  \draw[cRed, densely dotted, thick] (axis cs:2,4*pi) -- (axis cs:2.5,4*pi) -- (axis cs:2.5,5.9*pi);
  \node[cRed, below, font=\scriptsize] at (axis cs:2.25,4*pi) {$\Delta t$};
  \node[cRed, right, font=\scriptsize] at (axis cs:2.5,4.95*pi) {$\Delta\theta$ grande};
  \fill[cRed] (axis cs:2,4*pi) circle (1.5pt);
  \node[nota=cRed, anchor=west] at (axis cs:2.6,6) {$t_2$: inclinação grande,\\ $f_i(t_2)$ alta};
}
\only<5->{
  \addplot[cGreen, thick, domain=0.4:1.6] {2*pi + 0.2*pi*(x-1)};
  \draw[cGreen, densely dotted, thick] (axis cs:1,2*pi) -- (axis cs:1.5,2*pi) -- (axis cs:1.5,2.1*pi);
  \node[cGreen, below, font=\scriptsize] at (axis cs:1.25,2*pi) {$\Delta t$};
  \node[cGreen, below right, font=\scriptsize] at (axis cs:1.5,2.05*pi) {$\Delta\theta$ pequeno};
  \fill[cGreen] (axis cs:1,2*pi) circle (1.5pt);
  \node[nota=cGreen, anchor=south east] at (axis cs:0.95,10) {$t_1$: inclinação pequena,\\ $f_i(t_1)$ baixa};
}
\end{axis}
\end{tikzpicture}
\end{columns}

\only<3->{\anotacoes{
  \node[nota=cBlue, anchor=west] (n1) at ([xshift=0.45cm]fidtheta.east) {inclinação\\ de $\theta(t)$};
  \draw[setinha, cBlue] (n1.west) -- (fidtheta.east);
  \node[nota=cPurple, anchor=east] (n2) at ([xshift=-0.35cm]fi2pi.west) {rad/s $\to$ Hz};
  \draw[setinha, cPurple] (n2.east) -- (fi2pi.west);
}}

\end{frame}

% ============================================

\section{Modulação em Frequência (FM)}

\begin{frame}{Definição de FM}

\small
\begin{columns}[T]
\column{0.42\textwidth}
A frequência instantânea segue a mensagem:
\[
f_i(t) = f_c + \tikzmarknode{kf}{k_f}\, m(t)
\]

\vspace{0.6cm}
\onslide<3->{Desvio de frequência instantâneo
\[
\Delta f(t) = k_f\, m(t)
\]}
\onslide<4->{Desvio máximo
\[
\boxed{\Delta f = k_f \max|m(t)|}
\]}

\column{0.58\textwidth}
\begin{tikzpicture}
\begin{axis}[onda, width=6.6cm, height=1.3cm, domain=0:2, ymin=-1.2, ymax=1.2, ylabel={$m(t)$}]
\addplot[thick, cBlue] {0.7*sin(deg(2*pi*x)) + 0.3*sin(deg(6*pi*x))};
\only<4->{
  \addplot[cBlue, dashed, domain=0:2] {0.866};
  \node[font=\scriptsize, cBlue, anchor=south east] at (axis cs:2,0.87) {$\max|m|$};
}
\end{axis}
\onslide<2->{
\begin{axis}[onda, width=6.6cm, height=2.4cm, domain=0:2, ymin=-1.4, ymax=1.4, ylabel={$f_i(t)$}, at={(0,-3.1cm)}]
\addplot[gray, dashed] {0} node[pos=1, right, font=\scriptsize, gray] {$f_c$};
\addplot[very thick, cGreen] {0.7*sin(deg(2*pi*x)) + 0.3*sin(deg(6*pi*x))};
\only<4->{
  \fill[cGreen, opacity=0.12] (axis cs:0,-0.866) rectangle (axis cs:2,0.866);
  \draw[cRed, <->, thick] (axis cs:0.25,0) -- (axis cs:0.25,0.866) node[midway, right, font=\scriptsize] {$\Delta f$};
  \node[font=\scriptsize, cGreen!70!black, anchor=north east] at (axis cs:2,-0.9) {faixa $f_c \pm \Delta f$};
}
\end{axis}}
\end{tikzpicture}
\end{columns}

\only<2->{\anotacoes{
  \node[nota=cRed, anchor=north] (n1) at ([yshift=-0.35cm]kf.south) {Hz/V: mesma forma, outra escala};
  \draw[setinha, cRed] (n1.north) -- (kf.south);
}}

\end{frame}

% ============================================

\begin{frame}{Da frequência à fase: integração}

\small\compacto
Partimos da definição de frequência instantânea:
\[
\frac{d\theta(t)}{dt} = 2\pi f_i(t) = \tikzmarknode{c1}{2\pi f_c} + \tikzmarknode{c2}{2\pi k_f\, m(t)}
\]

\vspace{0.3cm}
\onslide<2->{Integramos os dois lados de $-\infty$ a $t$, termo a termo:
\[
\theta(t) = \tikzmarknode{i1}{\color{cBlue}2\pi f_c\, t} + \tikzmarknode{i2}{\color{cGreen}2\pi k_f \int_{-\infty}^{t} m(\tau)\,d\tau}
\]}

\vspace{0.2cm}
\onslide<4->{Comparando com $\theta(t) = 2\pi f_c t + \phi(t)$:
\[
\tikzmarknode{phibox}{\boxed{\;\phi(t) = 2\pi k_f \int_{-\infty}^{t} m(\tau)\,d\tau\;}}
\]}

\only<3->{\anotacoes{
  \node[nota=cBlue, anchor=east] (a1) at ([xshift=-1.6cm, yshift=-0.1cm]i1.west) {integral de constante:\\ cresce linear em $t$};
  \draw[setinha, cBlue] (a1.north east) to[bend left=35] (i1.north);
  \node[nota=cGreen, anchor=west] (a2) at ([xshift=0.5cm]i2.east) {isto é $\phi(t)$};
  \draw[setinha, cGreen] (a2.west) -- (i2.east);
  \draw[setinha, gray, dashed] (c1.south) -- (i1.north);
  \draw[setinha, gray, dashed] (c2.south) -- (i2.north);
}}
\only<5->{\anotacoes{
  \node[nota=cOrange, anchor=west] (a3) at ([xshift=0.9cm]phibox.east) {modulação com memória:\\ $\phi(t)$ depende de todo\\ o passado de $m(\tau)$, $\tau \le t$};
  \draw[setinha, cOrange] (a3.west) -- (phibox.east);
}}

\end{frame}

% ============================================

\begin{frame}{Exemplo: mensagem quadrada}

\small
\parbox[t][0.7cm]{\linewidth}{%
\only<1>{Mensagem alterna entre $+A$ e $-A$.}%
\only<2-5>{$\phi(t)$ acumula a área sob $m(t)$ até o instante $t$: área positiva faz a fase subir, área negativa faz a fase descer.}%
\only<6>{Integral de constante é rampa. Inclinação $\pm 2\pi k_f A$, logo $\phi(t)$ é triangular.}%
\only<7>{Rampas de fase deslocam a frequência: $f_c + k_f A$ nos trechos de subida, $f_c - k_f A$ nos de descida.}}

\begin{center}
\begin{tikzpicture}
\begin{axis}[onda, height=0.9cm, domain=0:4, xmin=0, xmax=4, ymin=-1.4, ymax=1.4, xlabel={$t$}, ylabel={$m(t)$}]
\only<2->{\fill[cGreen, opacity=0.25] (axis cs:0,0) rectangle (axis cs:1,1);}
\only<3->{\fill[cRed, opacity=0.25] (axis cs:1,0) rectangle (axis cs:2,-1);}
\only<4->{\fill[cGreen, opacity=0.25] (axis cs:2,0) rectangle (axis cs:3,1);}
\only<5->{\fill[cRed, opacity=0.25] (axis cs:3,0) rectangle (axis cs:4,-1);}
\addplot[thick, cBlue, const plot] coordinates {(0,1) (1,-1) (2,1) (3,-1) (4,-1)};
\node[font=\scriptsize, cBlue, anchor=south] at (axis cs:0.5,1) {$+A$};
\node[font=\scriptsize, cBlue, anchor=north] at (axis cs:1.5,-1) {$-A$};
\end{axis}
\onslide<2->{
\begin{axis}[onda, height=0.9cm, domain=0:4, xmin=0, xmax=4, ymin=-0.2, ymax=1.3, xlabel={$t$}, ylabel={$\phi(t)$}, at={(0,-1.75cm)}]
\only<2>{\addplot[very thick, cPurple] coordinates {(0,0) (1,1)};}
\only<3>{\addplot[very thick, cPurple] coordinates {(0,0) (1,1) (2,0)};}
\only<4>{\addplot[very thick, cPurple] coordinates {(0,0) (1,1) (2,0) (3,1)};}
\only<5->{\addplot[very thick, cPurple] coordinates {(0,0) (1,1) (2,0) (3,1) (4,0)};}
\only<2>{\fill[cPurple] (axis cs:1,1) circle (2pt);}
\only<3>{\fill[cPurple] (axis cs:2,0) circle (2pt);}
\only<4>{\fill[cPurple] (axis cs:3,1) circle (2pt);}
\only<5>{\fill[cPurple] (axis cs:4,0) circle (2pt);}
\only<6->{
  \node[font=\scriptsize, cPurple, anchor=south east] at (axis cs:0.5,0.55) {inclinação $+2\pi k_f A$};
  \node[font=\scriptsize, cPurple, anchor=south west] at (axis cs:1.5,0.55) {inclinação $-2\pi k_f A$};
}
\end{axis}}
\onslide<7->{
\begin{axis}[onda, height=1.15cm, domain=0:4, xmin=0, xmax=4, ymin=-1.3, ymax=1.3, xlabel={$t$}, ylabel={$s_{\FM}(t)$}, at={(0,-3.55cm)}, samples=300]
\addplot[cGreen, domain=0:1] {cos(deg(2*pi*6*x + 2*pi*3*x))};
\addplot[cGreen, domain=1:2] {cos(deg(2*pi*6*x + 2*pi*3*(2-x)))};
\addplot[cGreen, domain=2:3] {cos(deg(2*pi*6*x + 2*pi*3*(x-2)))};
\addplot[cGreen, domain=3:4] {cos(deg(2*pi*6*x + 2*pi*3*(4-x)))};
\end{axis}}
\end{tikzpicture}
\end{center}

\end{frame}

% ============================================

\begin{frame}{Sinal FM no domínio do tempo}

\small
Substituindo a fase na portadora:
\[
s_{\FM}(t) = \tikzmarknode{ac}{A_c} \cos\Big[2\pi f_c t + \tikzmarknode{beta}{2\pi k_f \int_{-\infty}^{t} m(\tau)\,d\tau}\Big]
\]

\vspace{0.3cm}
\onslide<3->{Chamando a fase modulada de $\beta(t)$:
\[
s_{\FM}(t) = A_c \cos[2\pi f_c t + \beta(t)]
\]}

\onslide<4->{
\begin{itemize}
\item Envoltória constante, $|s_{\FM}(t)| = A_c$: toda a informação está na fase.
\item Potência fixa, $P_s = A_c^2/2$, qualquer que seja $m(t)$.
\item Ruído ou desvanecimento que mexem na amplitude não tocam a mensagem.
\end{itemize}}

\only<2->{\anotacoes{
  \node[nota=cBlue, anchor=east] (n1) at ([xshift=-2.2cm]ac.west) {amplitude fixa};
  \draw[setinha, cBlue] (n1.east) to[bend left=25] (ac.north);
  \node[nota=cGreen, anchor=west] (n2) at ([xshift=0.7cm]beta.east) {$\beta(t)$: fase modulada};
  \draw[setinha, cGreen] (n2.west) to[bend right=25] (beta.north east);
}}

\end{frame}

% ============================================

\begin{frame}{Modulação em fase (PM)}

\small
\begin{columns}[T]
\column{0.45\textwidth}
A fase segue a mensagem, $k_p$ em rad/V:
\[
\phi(t) = k_p\, m(t)
\]
\[
s_{\PM}(t) = A_c \cos[2\pi f_c t + k_p m(t)]
\]

\vspace{0.1cm}
\onslide<2->{Frequência instantânea:
\[
f_i(t) = f_c + \frac{1}{2\pi}\frac{d\phi}{dt} = f_c + \frac{k_p}{2\pi}\,\tikzmarknode{dm}{\frac{dm(t)}{dt}}
\]}

\vspace{0.45cm}
\onslide<4->{Mensagem quadrada: fase salta, e $f_i(t)$ só muda nos instantes de salto.}

\column{0.55\textwidth}
\onslide<3->{
\begin{tikzpicture}
\begin{axis}[onda, width=6.6cm, height=1.0cm, domain=0:4, xmin=0, xmax=4, ymin=-1.4, ymax=1.4, ylabel={$m(t)$}]
\addplot[thick, cBlue, const plot] coordinates {(0,1) (1,-1) (2,1) (3,-1) (4,-1)};
\end{axis}
\begin{axis}[onda, width=6.6cm, height=2.0cm, domain=0:4, xmin=0, xmax=4, ymin=-1.3, ymax=1.3, ylabel={$s_{\PM}(t)$}, at={(0,-3.0cm)}, samples=300]
\addplot[cOrange, domain=0:1] {cos(deg(2*pi*6*x + pi/2))};
\addplot[cOrange, domain=1:2] {cos(deg(2*pi*6*x - pi/2))};
\addplot[cOrange, domain=2:3] {cos(deg(2*pi*6*x + pi/2))};
\addplot[cOrange, domain=3:4] {cos(deg(2*pi*6*x - pi/2))};
\only<4->{
  \foreach \k in {1,2,3} {\edef\tmp{\noexpand\draw[cRed, dashed] (axis cs:\k,-1.3) -- (axis cs:\k,1.3);}\tmp}
  \node[nota=cRed, anchor=north] at (axis cs:2,-1.35) {saltos de fase de $2k_pA$};
}
\end{axis}
\end{tikzpicture}}
\end{columns}

\only<2->{\anotacoes{
  \node[nota=cRed, anchor=north] (n1) at ([yshift=-0.5cm]dm.south) {derivada, não a própria $m$};
  \draw[setinha, cRed] (n1.north) -- (dm.south);
}}

\end{frame}

% ============================================

\begin{frame}{FM e PM com o mesmo tom}

\small
\parbox[t][0.75cm]{\linewidth}{%
\only<1>{Mensagem $m(t) = \cos(2\pi f_m t)$.}%
\only<2>{FM: $f_i \propto m(t)$. Oscilação mais rápida nos picos de $m$.}%
\only<3>{PM: $f_i \propto dm/dt$. Oscilação mais rápida onde $m$ mais cresce, um quarto de período depois.}%
\only<4>{Mesma forma de onda, deslocada no tempo: com tom único, FM e PM só se distinguem pela referência de fase.}}

\begin{center}
\begin{tikzpicture}
\begin{axis}[onda, height=0.9cm, domain=0:2, ymin=-1.3, ymax=1.3, ylabel={$m(t)$}]
\addplot[thick, cBlue] {cos(deg(2*pi*x))};
\only<2->{\foreach \k in {0,1,2} {\edef\tmp{\noexpand\fill[cGreen] (axis cs:\k,1) circle (1.6pt);}\tmp}}
\only<3->{\foreach \k in {0.75,1.75} {\edef\tmp{\noexpand\fill[cOrange] (axis cs:\k,0) circle (1.6pt);}\tmp}}
\end{axis}
\onslide<2->{
\begin{axis}[onda, height=1.2cm, domain=0:2, ymin=-1.3, ymax=1.3, ylabel={FM}, at={(0,-1.95cm)}, samples=800]
\addplot[cGreen] {cos(deg(2*pi*8*x + 5*sin(deg(2*pi*x))))};
\foreach \k in {0,1,2} {\edef\tmp{\noexpand\draw[cGreen, dashed] (axis cs:\k,-1.3) -- (axis cs:\k,1.3);}\tmp}
\end{axis}}
\onslide<3->{
\begin{axis}[onda, height=1.1cm, domain=0:2, ymin=-1.3, ymax=1.3, ylabel={PM}, at={(0,-3.55cm)}, samples=800]
\addplot[cOrange] {cos(deg(2*pi*8*x + 5*cos(deg(2*pi*x))))};
\foreach \k in {0.75,1.75} {\edef\tmp{\noexpand\draw[cOrange, dashed] (axis cs:\k,-1.3) -- (axis cs:\k,1.3);}\tmp}
\end{axis}}
\end{tikzpicture}
\end{center}

\end{frame}

% ============================================

\begin{frame}{FM versus PM: comparação direta}

\begin{center}
\renewcommand{\arraystretch}{1.5}
\begin{tabular}{lcc}
\toprule
                        & FM                                 & PM \\
\midrule
Fase $\phi(t)$          & $2\pi k_f \int m(\tau)\,d\tau$     & $k_p\, m(t)$ \\
\onslide<2->{Frequência $f_i(t) - f_c$} & \onslide<2->{$k_f\, m(t)$} & \onslide<2->{$\dfrac{k_p}{2\pi}\,\dfrac{dm(t)}{dt}$} \\
\bottomrule
\end{tabular}
\end{center}

\vspace{0.3cm}
\begin{itemize}
\item<3-> Em FM, a frequência copia a mensagem; em PM, é a fase que copia.
\item<4-> FM de $m(t)$ é PM de $\int m\,dt$.
\item<5-> PM de $m(t)$ é FM de $dm/dt$.
\item<6-> Na prática, FM é mais usada por ter melhor desempenho em ruído.
\end{itemize}

\end{frame}

% ============================================

\begin{frame}{Conversão prática FM--PM}

\begin{center}
\begin{tikzpicture}[>=latex]
\node[draw, rectangle, minimum width=1.5cm, minimum height=0.7cm] (int) at (0,1.2) {$\int dt$};
\node[draw, rectangle, minimum width=1.5cm, minimum height=0.7cm] (pm) at (3,1.2) {Mod.\ PM};
\draw[->] (-1.7,1.2) -- node[above, font=\small] {$m(t)$} (int);
\draw[->] (int) -- (pm);
\draw[->] (pm) -- (5,1.2) node[right, font=\small] {FM de $m(t)$};

\onslide<2->{
\node[draw, rectangle, minimum width=1.5cm, minimum height=0.7cm] (der) at (0,-1.2) {$d/dt$};
\node[draw, rectangle, minimum width=1.5cm, minimum height=0.7cm] (fm) at (3,-1.2) {Mod.\ FM};
\draw[->] (-1.7,-1.2) -- node[above, font=\small] {$m(t)$} (der);
\draw[->] (der) -- (fm);
\draw[->] (fm) -- (5,-1.2) node[right, font=\small] {PM de $m(t)$};}
\end{tikzpicture}
\end{center}

\vspace{0.3cm}
\small
\onslide<3->{Um modulador PM precedido de integrador produz FM. É o caminho do método de Armstrong, que veremos adiante.}

\end{frame}

% ============================================

\section{Análise Espectral de FM}

\begin{frame}{FM com tom único}

\small
Mensagem com fase qualquer, $m(t) = A_m \cos(2\pi f_m t + \theta_m)$:
\[
s(t) = A_c \cos\Big[2\pi f_c t + 2\pi k_f A_m \tikzmarknode{int}{\int_{-\infty}^{t} \cos(2\pi f_m \tau + \theta_m)\,d\tau}\Big]
\]

\vspace{0.35cm}
\onslide<2->{Integral do cosseno:
\[
\tikzmarknode{lhs}{\int_{-\infty}^{t} \cos(2\pi f_m \tau + \theta_m)\,d\tau} = \frac{\sin(2\pi f_m t + \theta_m)}{\tikzmarknode{den}{2\pi f_m}} + \tikzmarknode{cte}{\color{cBlue}C}
\]}

\vspace{0.1cm}
\onslide<3->{Substituindo:
\[
s(t) = A_c \cos\Big[2\pi f_c t + \frac{\cancel{2\pi}\, k_f A_m}{\cancel{2\pi} f_m}\,\sin(2\pi f_m t + \theta_m)\Big]
     = A_c \cos\Big[2\pi f_c t + \tikzmarknode{res}{\color{cGreen}\frac{k_f A_m}{f_m}}\,\sin(2\pi f_m t + \theta_m)\Big]
\]}

\only<2->{\anotacoes{
  \node[nota=cBlue, anchor=south east] (n1) at ([xshift=0.4cm, yshift=0.45cm]cte.north east) {$C$: ângulo fixo somado à portadora,\\ não leva $m(t)$. Escolhemos $C = 0$};
  \draw[setinha, cBlue] (n1.south east) -- (cte.north);
}}
\only<3>{\anotacoes{
  \node[nota=cRed, anchor=east] (n2) at ([xshift=-0.4cm]lhs.west) {derivada do seno:\\ $2\pi f_m\cos(\cdot)$};
  \draw[setinha, cRed] (n2.south) to[bend right=25] (den.south west);
}}
\only<4->{\anotacoes{
  \node[nota=cGreen, anchor=south] (n3) at ([yshift=0.45cm]res.north) {adimensional: chamaremos de $\beta$};
  \draw[setinha, cGreen] (n3.south) -- (res.north);
}}

\end{frame}

% ============================================

\begin{frame}{Índice de modulação $\beta$}

\small
\begin{columns}[T]
\column{0.44\textwidth}
\[
\boxed{\;\beta = \frac{k_f A_m}{f_m} = \frac{\tikzmarknode{df}{\Delta f}}{\tikzmarknode{fm}{f_m}}\;}
\]
\vspace{0.6cm}
\[
s(t) = A_c \cos[2\pi f_c t + \beta \sin(2\pi f_m t + \theta_m)]
\]

\vspace{-0.1cm}
\onslide<2->{Leitura 1: $\beta$ é o pico do desvio de fase, em radianos.}

\vspace{0.02cm}
\onslide<3->{Leitura 2: $\beta$ compara o desvio de frequência com a frequência da mensagem.}

\vspace{0.02cm}
\onslide<4->{Mesmo $\Delta f$, $f_m$ pela metade: o desvio dura o dobro do tempo, a fase acumula o dobro e $\beta$ dobra.}

\column{0.56\textwidth}
\onslide<2->{
\begin{tikzpicture}
\begin{axis}[onda, width=6.4cm, height=1.6cm, domain=0:2, ymin=-4.4, ymax=4.4, ylabel={$\phi(t)$}]
\only<2-3>{
  \addplot[very thick, cPurple] {2*sin(deg(2*pi*x))};
  \draw[cPurple, <->] (axis cs:0.25,0) -- (axis cs:0.25,2) node[midway, right, font=\scriptsize] {$\beta$};
}
\only<4->{
  \addplot[gray, dashed] {2*sin(deg(2*pi*x))};
  \addplot[very thick, cPurple] {4*sin(deg(pi*x))};
  \draw[cPurple, <->] (axis cs:0.5,0) -- (axis cs:0.5,4) node[midway, right, font=\scriptsize] {$\beta$ dobra};
}
\end{axis}
\onslide<3->{
\begin{axis}[onda, width=6.4cm, height=1.6cm, domain=0:2, ymin=-1.3, ymax=1.3, ylabel={$f_i(t) - f_c$}, at={(0,-2.5cm)}]
\only<3>{\addplot[very thick, cGreen] {cos(deg(2*pi*x))};}
\only<4->{\addplot[gray, dashed] {cos(deg(2*pi*x))}; \addplot[very thick, cGreen] {cos(deg(pi*x))};}
\draw[cGreen, <->] (axis cs:1.6,0) -- (axis cs:1.6,1) node[midway, right, font=\scriptsize] {$\Delta f$ igual};
\end{axis}}
\end{tikzpicture}}
\end{columns}

\only<2->{\anotacoes{
  \node[nota=cGreen, anchor=north] (n1) at ([xshift=-0.6cm, yshift=-0.3cm]fm.south) {desvio de pico / ritmo da mensagem};
  \draw[setinha, cGreen] (n1.north) -- (df.south);
}}

\end{frame}

% ============================================

\begin{frame}{A fase da mensagem só desloca o tempo}

\small
\centering
\begin{tikzpicture}
\pgfplotsset{rot/.style={ylabel={}, width=11.5cm}}
\begin{axis}[onda, rot, height=0.6cm, domain=0:2, ymin=-1.3, ymax=1.3]
\node[anchor=east, font=\scriptsize] at (axis cs:-0.03,0) {$m(t)$};
\addplot[thick, cBlue, dashed] {cos(deg(2*pi*x))};
\addplot[thick, cRed] {cos(deg(2*pi*x + pi/2))};
\node[anchor=south east, font=\scriptsize, fill=white, inner sep=1pt] at (axis cs:2,1.3) {\textcolor{cBlue}{tracejado: $\theta_m = 0$}, \textcolor{cRed}{cheio: $\theta_m = \pi/2$}};
\only<3->{
  \draw[cBlue, dashed, thin] (axis cs:1,-1.3) -- (axis cs:1,1.3);
  \draw[cRed, dashed, thin] (axis cs:0.75,-1.3) -- (axis cs:0.75,1.3);
  \draw[cPurple, <->, thick] (axis cs:0.75,1.6) -- (axis cs:1,1.6) node[midway, above, font=\scriptsize] {$t_0$};
}
\end{axis}
\onslide<2->{
\begin{axis}[onda, rot, height=0.7cm, domain=0:2, ymin=-1.3, ymax=1.3, at={(0,-1.05cm)}, samples=800]
\node[anchor=east, font=\scriptsize, cBlue] at (axis cs:-0.03,0) {$\theta_m = 0$};
\addplot[cBlue] {cos(deg(2*pi*8*x + 5*sin(deg(2*pi*x))))};
\only<3->{\draw[cBlue, dashed, thin] (axis cs:1,-1.3) -- (axis cs:1,1.3);}
\end{axis}
\begin{axis}[onda, rot, height=0.7cm, domain=0:2, ymin=-1.3, ymax=1.3, at={(0,-2.0cm)}, samples=800]
\node[anchor=east, font=\scriptsize, cRed] at (axis cs:-0.03,0) {$\theta_m = \pi/2$};
\addplot[cRed] {cos(deg(2*pi*8*x + 5*sin(deg(2*pi*x + pi/2))))};
\only<3->{\draw[cRed, dashed, thin] (axis cs:0.75,-1.3) -- (axis cs:0.75,1.3);}
\end{axis}}
\end{tikzpicture}

\raggedright
\onslide<3->{Com $t_0 = \theta_m/(2\pi f_m)$, o FM com $\theta_m$ é o FM com $\theta_m = 0$ adiantado de $t_0$:
\[
2\pi f_c t + \beta\sin(2\pi f_m t + \theta_m) = 2\pi f_c (t + t_0) + \beta\sin\big(2\pi f_m (t + t_0)\big) - \tikzmarknode{fase}{2\pi f_c t_0}
\]}

\vspace{0.1cm}
\onslide<4->{Vamos carregar $\theta_m$ nas contas e ver onde ela aparece no espectro.}

\only<3->{\anotacoes{
  \node[nota=cPurple, anchor=west] (n1) at ([xshift=0.3cm, yshift=-0.45cm]fase.east) {fase constante};
  \draw[setinha, cPurple] (n1.west) -- (fase.south east);
}}

\end{frame}

% ============================================

\begin{frame}{Bessel, passo 1: o FM em banda básica}

\small\compacto
\begin{columns}[T]
\column{0.6\textwidth}
Com $\omega_c = 2\pi f_c$ e $\omega_m = 2\pi f_m$:
\[
s(t) = A_c \cos[\omega_c t + \beta\sin(\omega_m t + \theta_m)]
\]

\onslide<2->{Como $\cos x = \Re\{e^{jx}\}$:
\[
s(t) = \Re\Big\{A_c\, \tikzmarknode{port}{e^{j\omega_c t}}\; \tikzmarknode{g}{e^{j\beta\sin(\omega_m t + \theta_m)}}\Big\}
\]}

\vspace{0.6cm}
\onslide<3->{Envoltória complexa, em banda básica:
\[
\boxed{\,g(t) = e^{j\beta\sin(\omega_m t + \theta_m)}\,}
\]}

\column{0.4\textwidth}
\onslide<3->{
\centering
\begin{tikzpicture}[scale=1.1]
\draw[->, gray] (-1.3,0) -- (1.45,0) node[right, font=\tiny] {$\Re$};
\draw[->, gray] (0,-1.3) -- (0,1.4) node[above, font=\tiny] {$\Im$};
\draw[gray!60] (0,0) circle (1);
\draw[cPurple, line width=2.2pt, opacity=0.35] (-63:1) arc (-63:63:1);
\node[cPurple, font=\scriptsize] at (70:1.25) {$+\beta$};
\node[cPurple, font=\scriptsize] at (-70:1.25) {$-\beta$};
\only<4>{\draw[-{Stealth}, very thick, cGreen] (0,0) -- (0:1); \node[font=\scriptsize, cGreen, anchor=west] at (0:1.02) {$t = 0$};}
\only<5>{\draw[-{Stealth}, very thick, cGreen] (0,0) -- (63:1); \node[font=\scriptsize, cGreen, anchor=south west] at (63:1.02) {$t = T_m/4$};}
\only<6>{\draw[-{Stealth}, very thick, cGreen] (0,0) -- (0:1); \node[font=\scriptsize, cGreen, anchor=west] at (0:1.02) {$t = T_m/2$};}
\only<7>{\draw[-{Stealth}, very thick, cGreen] (0,0) -- (-63:1); \node[font=\scriptsize, cGreen, anchor=north west] at (-63:1.02) {$t = 3T_m/4$};}
\draw[gray, ->] (0.35,0) arc (0:63:0.35);
\node[font=\tiny, gray] at (30:0.55) {$\beta$};
\end{tikzpicture}

{\scriptsize $g(t)$ vai e volta no círculo unitário, entre $-\beta$ e $+\beta$.}

\onslide<7->{\scriptsize $g(t + T_m) = g(t)$, $T_m = 2\pi/\omega_m$: periódica, admite série de Fourier.}}
\end{columns}

\only<2->{\anotacoes{
  \node[nota=cBlue, anchor=north] (a1) at ([xshift=-0.6cm, yshift=-0.3cm]port.south) {portadora};
  \draw[setinha, cBlue] (a1.north) -- (port.south);
  \node[nota=cGreen, anchor=north] (a2) at ([yshift=-0.3cm]g.south) {só depende da mensagem};
  \draw[setinha, cGreen] (a2.north) -- (g.south);
}}

\end{frame}

% ============================================

\begin{frame}{Bessel, passo 2: série de Fourier}

\small
\begin{columns}[T]
\column{0.5\textwidth}
Toda função periódica de período $T_m$ admite
\[
g(t) = \sum_{n=-\infty}^{\infty} c_n\, e^{jn\omega_m t}
\]

\onslide<3->{com coeficientes
\[
c_n = \frac{1}{T_m}\int_{0}^{T_m} \tikzmarknode{gg}{e^{j\beta\sin(\omega_m t + \theta_m)}}\;\tikzmarknode{ee}{e^{-jn\omega_m t}}\,dt
\]}

\vspace{0.6cm}
\onslide<4->{Falta calcular essa integral. É o que faremos agora, sem pular passos.}

\column{0.5\textwidth}
\begin{tikzpicture}
\begin{axis}[onda, width=6.2cm, height=1.3cm, domain=0:2, ymin=-1.3, ymax=1.3, ylabel={$\Re\{g(t)\}$}]
\addplot[very thick, cBlue] {cos(deg(2*sin(deg(2*pi*x))))};
\onslide<2->{
  \draw[cRed, |<->|] (axis cs:0,-1.25) -- (axis cs:1,-1.25) node[midway, below, font=\scriptsize] {$T_m$};
  \draw[cRed, |<->|] (axis cs:1,-1.25) -- (axis cs:2,-1.25) node[midway, below, font=\scriptsize] {$T_m$};}
\end{axis}
\begin{axis}[onda, width=6.2cm, height=1.3cm, domain=0:2, ymin=-1.3, ymax=1.3, ylabel={$\Im\{g(t)\}$}, at={(0,-2.3cm)}]
\addplot[very thick, cOrange] {sin(deg(2*sin(deg(2*pi*x))))};
\end{axis}
\end{tikzpicture}

{\scriptsize $\beta = 2$, $\theta_m = 0$: cada período se repete idêntico.}
\end{columns}

\only<3->{\anotacoes{
  \node[nota=cGreen, anchor=north] (a1) at ([xshift=-0.5cm, yshift=-0.35cm]gg.south) {$g(t)$};
  \draw[setinha, cGreen] (a1.north) -- (gg.south);
  \node[nota=cBlue, anchor=north west] (a2) at ([xshift=0.1cm, yshift=-0.35cm]ee.south) {projeção na harmônica $n$};
  \draw[setinha, cBlue] (a2.north west) -- (ee.south);
}}

\end{frame}

% ============================================

\begin{frame}{Bessel, passo 3: mudança de variável}

\small
\begin{columns}[T]
\column{0.64\textwidth}
\[
c_n = \frac{1}{T_m}\int_{0}^{T_m} e^{j\beta\sin(\tikzmarknode{arg}{\omega_m t + \theta_m})}\; e^{-jn\omega_m t}\,dt
\]

\vspace{0.2cm}
\onslide<2->{Troca $x = \omega_m t + \theta_m$:\quad $dt = \dfrac{dx}{\omega_m}$,\quad $\omega_m t = x - \theta_m$.}

\vspace{0.25cm}
\onslide<3->{\[
c_n = \frac{1}{\tikzmarknode{wT}{\omega_m T_m}}\int_{\tikzmarknode{li}{\theta_m}}^{\theta_m + 2\pi} e^{j\beta\sin x}\; \tikzmarknode{ex}{e^{-jn(x - \theta_m)}}\,dx
\]}

\vspace{0.4cm}
\onslide<5->{\[
c_n = \tikzmarknode{fora}{\color{cPurple}e^{jn\theta_m}}\;\frac{1}{2\pi}\int_{\theta_m}^{\theta_m + 2\pi} e^{j(\beta\sin x - nx)}\,dx
\]}

\column{0.36\textwidth}
\onslide<2->{
\begin{tikzpicture}
\begin{axis}[espec, width=3.6cm, height=2.9cm, xmin=0, xmax=1.25, ymin=0, ymax=8.2,
             xtick={0,1}, xticklabels={$0$, $T_m$}, ytick={1,7.28}, yticklabels={$\theta_m$, $\theta_m + 2\pi$},
             xlabel={$t$}, ylabel={$x$}]
\addplot[very thick, cGreen, domain=0:1] {1 + 2*pi*x};
\draw[dashed, gray] (axis cs:1,0) -- (axis cs:1,7.28) -- (axis cs:0,7.28);
\fill[cGreen] (axis cs:0,1) circle (1.6pt);
\fill[cGreen] (axis cs:1,7.28) circle (1.6pt);
\node[font=\tiny, cGreen, anchor=north west, align=left] at (axis cs:0.35,3) {inclinação\\ $\omega_m$};
\end{axis}
\end{tikzpicture}

{\scriptsize $t: 0 \to T_m$ vira $x: \theta_m \to \theta_m + 2\pi$.}}
\end{columns}

\only<3-4>{\anotacoes{
  \node[nota=cRed, anchor=north] (a1) at ([yshift=-0.35cm]wT.south) {$\omega_m T_m = 2\pi$};
  \draw[setinha, cRed] (a1.north) -- (wT.south);
}}
\only<4>{\anotacoes{
  \node[nota=cPurple, anchor=north] (a3) at ([yshift=-0.35cm]ex.south) {$= e^{jn\theta_m}\, e^{-jnx}$: constante sai};
  \draw[setinha, cPurple] (a3.north) -- (ex.south);
}}

\end{frame}

% ============================================

\begin{frame}{Bessel, passo 4: janela de $2\pi$}

\small\compacto
O integrando $e^{j(\beta\sin x - nx)}$ tem período $2\pi$ em $x$. Parte real, com $\beta = 2$ e $n = 1$:

\vspace{0.1cm}
\centering
\begin{tikzpicture}
\begin{axis}[onda, width=11cm, height=1.1cm, domain=-3.6:7.6, samples=300, ymin=-1.3, ymax=1.3,
             xlabel={$x$}, ylabel={}]
\only<2->{\fill[cGreen, opacity=0.15] (axis cs:1,-1.3) rectangle (axis cs:7.283,1.3);
  \node[font=\scriptsize, cGreen, anchor=south] at (axis cs:4.14,1.3) {$[\theta_m,\ \theta_m + 2\pi]$};}
\only<3->{\fill[cBlue, opacity=0.15] (axis cs:-3.1416,-1.3) rectangle (axis cs:3.1416,1.3);
  \node[font=\scriptsize, cBlue, anchor=south] at (axis cs:-1.2,1.3) {$[-\pi,\ \pi]$};}
\only<4->{\draw[cRed, ->, thick] (axis cs:5.4,-1.5) to[bend left=12] node[midway, below, font=\scriptsize, align=center] {$[\pi,\ \theta_m + 2\pi]$ é cópia de $[-\pi,\ \theta_m]$: mesma área} (axis cs:-0.9,-1.5);}
\addplot[thick, cOrange] {cos(deg(2*sin(deg(x)) - x))};
\end{axis}
\end{tikzpicture}

\vspace{0.15cm}
\raggedright
\onslide<5->{\[
c_n = \tikzmarknode{efase}{e^{jn\theta_m}}\;\underbrace{\frac{1}{2\pi}\int_{-\pi}^{\pi} e^{j(\beta\sin x - nx)}\,dx}_{\textstyle \triangleq\, J_n(\beta)\text{, Bessel de ordem } n}
\qquad\Longrightarrow\qquad
\boxed{\,c_n = J_n(\beta)\, e^{jn\theta_m}\,}
\]}

\only<5->{\anotacoes{
  \node[nota=cPurple, anchor=east] (a1) at ([xshift=-0.6cm]efase.west) {$\theta_m$ só entra aqui};
  \draw[setinha, cPurple] (a1.east) -- (efase.west);
}}

\end{frame}

% ============================================

\begin{frame}{Funções de Bessel $J_n(\beta)$}

\small
\begin{columns}[T]
\column{0.6\textwidth}
\begin{tikzpicture}
\begin{axis}[espec, width=7cm, height=3.6cm, xmin=0, xmax=10, ymin=-0.45, ymax=1.05,
             axis lines=middle, xlabel={$\beta$}, ylabel={$J_n(\beta)$}, xtick={0,2,4,6,8,10}, ytick={-0.4,0,0.5,1}]
\addplot[very thick, cBlue] table[x=beta, y=J0]{dados/bessel_curvas.dat} node[pos=0.03, above right, font=\scriptsize] {$J_0$};
\onslide<2->{\addplot[very thick, cRed] table[x=beta, y=J1]{dados/bessel_curvas.dat};
  \node[font=\scriptsize, cRed] at (axis cs:1.8,0.66) {$J_1$};}
\onslide<3->{\addplot[very thick, cGreen] table[x=beta, y=J2]{dados/bessel_curvas.dat};
  \node[font=\scriptsize, cGreen] at (axis cs:3.1,0.55) {$J_2$};
  \addplot[very thick, cOrange] table[x=beta, y=J3]{dados/bessel_curvas.dat};
  \node[font=\scriptsize, cOrange] at (axis cs:4.3,0.5) {$J_3$};}
\only<4->{\fill[cPurple] (axis cs:2.405,0) circle (2pt);
  \node[nota=cPurple, anchor=north west] at (axis cs:2.5,-0.05) {$\beta = 2{,}405$: $J_0 = 0$};}
\end{axis}
\end{tikzpicture}

\column{0.4\textwidth}
\onslide<1->{$J_0(\beta)$ pesa a portadora.}

\vspace{0.15cm}
\onslide<2->{$J_n$ só cresce depois de $\beta \approx n$: com $\beta$ maior, mais componentes acordam.}

\vspace{0.15cm}
\onslide<3->{$\beta$ pequeno: $J_n(\beta) \approx \dfrac{(\beta/2)^n}{n!}$.}

\vspace{0.15cm}
\onslide<4->{Em $\beta = 2{,}405$ a portadora some: toda a potência vai para as laterais.}

\vspace{0.15cm}
\onslide<5->{Potência: $\displaystyle\sum_{n} J_n^2(\beta) = 1$.}
\end{columns}

\end{frame}

% ============================================

\begin{frame}{Simetria de $J_n(\beta)$}

\small\compacto
\begin{columns}[T]
\column{0.56\textwidth}
\[
J_{-n}(\beta) = \frac{1}{2\pi}\int_{-\pi}^{\pi} e^{j(\beta\sin x \tikzmarknode{mais}{+ nx})}\,dx
\]

\vspace{0.3cm}
\onslide<2->{Troca $x = \pi - u$: \ $\sin x = \sin u$, \ $nx = n\pi - nu$, \ $dx = -du$; o $-du$ inverte os limites.}

\vspace{0.2cm}
\onslide<3->{\[
J_{-n}(\beta) = \tikzmarknode{epi}{e^{jn\pi}}\;\frac{1}{2\pi}\int_{0}^{2\pi} e^{j(\beta\sin u - nu)}\,du
\]}

\vspace{0.3cm}
\onslide<4->{\[
\boxed{\,J_{-n}(\beta) = (-1)^n J_n(\beta)\,}
\]
Ordem ímpar troca de sinal; o módulo é o mesmo.}

\column{0.44\textwidth}
\begin{tikzpicture}
\begin{axis}[espec, width=5.2cm, height=3.4cm, axis lines=middle, xmin=-5.8, xmax=5.8, ymin=-0.7, ymax=0.7,
             xlabel={$n$}, ylabel={$J_n(2)$}, xtick={-5,-3,-1,1,3,5}, ytick={-0.5,0.5}]
\addplot+[comp, cBlue, mark options={fill=cBlue}] table[x=n, y=J, restrict x to domain=0:6]{dados/comp_b2.dat};
\onslide<4->{\addplot+[comp, cRed, mark options={fill=cRed}] table[x=n, y=J, restrict x to domain=-6:-1]{dados/comp_b2.dat};}
\end{axis}
\end{tikzpicture}
\end{columns}

\only<3->{\anotacoes{
  \node[nota=cRed, anchor=north] (a1) at ([yshift=-0.35cm]epi.south) {$e^{jn\pi} = (-1)^n$};
  \draw[setinha, cRed] (a1.north) -- (epi.south);
}}

\end{frame}

% ============================================

\begin{frame}{Espectro do FM com tom único}

\small\compacto
\begin{columns}[T]
\column{0.58\textwidth}
\[
g(t) = \sum_{n} c_n e^{jn\omega_m t} = \sum_{n} J_n(\beta)\, e^{jn(\omega_m t + \theta_m)}
\]

\onslide<2->{\[
s(t) = \Re\{A_c\, e^{j\omega_c t}\, g(t)\}
\]}

\onslide<3->{\[
\boxed{\,s(t) = A_c\sum_{n=-\infty}^{\infty} \tikzmarknode{amp}{J_n(\beta)}\,\cos[(\omega_c + n\omega_m)t + \tikzmarknode{fs}{n\theta_m}]\,}
\]}


\column{0.42\textwidth}
\onslide<4->{
\begin{tikzpicture}
\begin{axis}[espec, width=4.8cm, height=3.0cm, xmin=-6, xmax=6, ymax=0.7,
             xlabel={$f$}, ylabel={$A_c|J_n(2)|$}, xtick={-4,-2,0,2,4},
             xticklabels={$f_c{-}4f_m$, , $f_c$, , $f_c{+}4f_m$}, ytick={0.2,0.4,0.6}]
\addplot+[comp, cBlue, mark options={fill=cBlue}] table[x=n, y=absJ]{dados/comp_b2.dat};
\end{axis}
\end{tikzpicture}

{\scriptsize $\beta = 2$: componentes de frequência em $f_c + nf_m$.}}
\end{columns}

\only<3->{\anotacoes{
  \node[nota=cBlue, anchor=north] (a1) at ([xshift=-0.6cm, yshift=-0.3cm]amp.south) {amplitude: não depende de $\theta_m$};
  \draw[setinha, cBlue] (a1.north) -- (amp.south);
  \node[nota=cPurple, anchor=north] (a2) at ([xshift=0.5cm, yshift=-0.3cm]fs.south) {$\theta_m$ só gira a fase};
  \draw[setinha, cPurple] (a2.north) -- (fs.south);
}}

\end{frame}

% ============================================

\begin{frame}{O espectro conforme $\beta$ cresce}

\small
\parbox[t][0.55cm]{\linewidth}{%
\only<1>{$\beta = 0{,}5$: portadora dominante e um par de laterais, como em AM.}%
\only<2>{$\beta = 1$: o segundo par aparece.}%
\only<3>{$\beta = 2$: a portadora já perdeu a maior parte da potência.}%
\only<4>{$\beta = 5$: muitas componentes, portadora pequena. A faixa sombreada é a regra de Carson, $|n| \le \beta + 1$.}}

\begin{center}
\begin{tikzpicture}
\begin{axis}[espec, width=11cm, height=3.4cm, xmin=-9.5, xmax=9.5, ymax=1.0,
             xlabel={$n$}, ylabel={$|J_n(\beta)|$}, xtick={-8,-6,...,8}, ytick={0.25,0.5,0.75,1}]
\only<1>{\fill[cGreen, opacity=0.12] (axis cs:-1.5,0) rectangle (axis cs:1.5,1);
  \addplot+[comp, cBlue, mark options={fill=cBlue}] table[x=n, y=absJ]{dados/comp_b0p5.dat};}
\only<2>{\fill[cGreen, opacity=0.12] (axis cs:-2.5,0) rectangle (axis cs:2.5,1);
  \addplot+[comp, cBlue, mark options={fill=cBlue}] table[x=n, y=absJ]{dados/comp_b1.dat};}
\only<3>{\fill[cGreen, opacity=0.12] (axis cs:-3.5,0) rectangle (axis cs:3.5,1);
  \addplot+[comp, cBlue, mark options={fill=cBlue}] table[x=n, y=absJ]{dados/comp_b2.dat};}
\only<4>{\fill[cGreen, opacity=0.12] (axis cs:-6.5,0) rectangle (axis cs:6.5,1);
  \addplot+[comp, cBlue, mark options={fill=cBlue}] table[x=n, y=absJ]{dados/comp_b5.dat};}
\end{axis}
\end{tikzpicture}
\end{center}

\only<4>{Na faixa de Carson ficam 99{,}4\,\% da potência para $\beta = 5$.}

\end{frame}

% ============================================

\begin{frame}{Exercício: $\beta = 2$}

\small
$A_c = 1$\,V, $f_c = 100$\,kHz, $f_m = 1$\,kHz, $\beta = 2$. Quais componentes têm $|J_n| \ge 0{,}01$?

\begin{columns}[T]
\column{0.42\textwidth}
\onslide<2->{
\scriptsize
\renewcommand{\arraystretch}{1.05}
\begin{tabular}{ccc}
\toprule
$n$ & $f$ [kHz] & $|J_n(2)|$ \\
\midrule
$0$       & 100      & 0,224 \\
$\pm 1$   & 99, 101  & 0,577 \\
$\pm 2$   & 98, 102  & 0,353 \\
$\pm 3$   & 97, 103  & 0,129 \\
$\pm 4$   & 96, 104  & 0,034 \\
$\pm 5$   & 95, 105  & 0,007 \\
\bottomrule
\end{tabular}}

\vspace{0.2cm}
\onslide<3->{$|n| \le 4$: nove componentes, de 96 a 104 kHz.}

\column{0.58\textwidth}
\onslide<2->{
\begin{tikzpicture}
\begin{axis}[espec, width=6cm, height=2.9cm, xmin=94, xmax=106, ymax=0.7,
             xlabel={$f$ [kHz]}, ylabel={amplitude [V]}, xtick={95,97,...,105}, ytick={0.2,0.4,0.6}]
\addplot+[comp, cBlue, mark options={fill=cBlue}] table[x expr=\thisrow{n}+100, y=absJ]{dados/comp_b2.dat};
\only<3->{\addplot[cRed, dashed, domain=94:106] {0.01} node[pos=0, above right, font=\tiny] {0,01};}
\end{axis}
\end{tikzpicture}}

\onslide<4->{Potência: $J_0^2 + 2\sum_{n=1}^{4} J_n^2(2) = 0{,}050 + 2(0{,}475) = 0{,}9999$.}
\end{columns}

\end{frame}

% ============================================

\begin{frame}{Dois tons: o sinal}

\small\compacto
\begin{columns}[T]
\column{0.5\textwidth}
\[
m(t) = A_1\cos(\omega_1 t + \theta_1) + A_2\cos(\omega_2 t + \theta_2)
\]

\onslide<2->{A fase é a soma de duas fases de tom único:
\[
\phi(t) = \beta_1\sin(\omega_1 t + \theta_1) + \beta_2\sin(\omega_2 t + \theta_2)
\]}

\onslide<3->{Em banda básica, exponencial de soma vira produto:
\[
g(t) = \tikzmarknode{g1}{e^{j\beta_1\sin(\omega_1 t + \theta_1)}}\;\tikzmarknode{g2}{e^{j\beta_2\sin(\omega_2 t + \theta_2)}}
\]}

\vspace{0.5cm}
\onslide<4->{\scriptsize Exemplo: $f_1 = 1$\,kHz, $f_2 = 10$\,kHz, $\beta_1 = 1$, $\beta_2 = 0{,}5$, $A_c = 1$\,V, $f_c = 100$\,kHz.}

\column{0.5\textwidth}
\begin{tikzpicture}
\begin{axis}[onda, width=5.8cm, height=1.2cm, domain=0:1, samples=400, ymin=-1.6, ymax=1.6, ylabel={$m(t)$}]
\addplot[thick, cBlue] {0.6*cos(deg(2*pi*2*x)) + 0.4*cos(deg(2*pi*15*x + 1))};
\end{axis}
\onslide<2->{
\begin{axis}[onda, width=5.8cm, height=1.2cm, domain=0:1, samples=400, ymin=-1.8, ymax=1.8, ylabel={$\phi(t)$}, at={(0,-2.1cm)}]
\addplot[thick, cPurple] {1.2*sin(deg(2*pi*2*x)) + 0.5*sin(deg(2*pi*15*x + 1))};
\addplot[cPurple!50, dashed] {1.2*sin(deg(2*pi*2*x))};
\end{axis}}
\end{tikzpicture}
\end{columns}

\only<3->{\anotacoes{
  \node[nota=cBlue, anchor=north] (a1) at ([yshift=-0.3cm]g1.south) {$g_1(t)$};
  \draw[setinha, cBlue] (a1.north) -- (g1.south);
  \node[nota=cRed, anchor=north] (a2) at ([yshift=-0.3cm]g2.south) {$g_2(t)$};
  \draw[setinha, cRed] (a2.north) -- (g2.south);
}}

\end{frame}

% ============================================

\begin{frame}{Dois tons: convolução}

\small
\parbox[t][1.15cm]{\linewidth}{%
\only<1>{Espectro de $g_1(t)$: componentes $J_m(\beta_1)\,e^{jm\theta_1}$ em $m f_1$, espaçadas de 1 kHz.}%
\only<2>{Espectro de $g_2(t)$: componentes $J_n(\beta_2)\,e^{jn\theta_2}$ em $n f_2$, espaçadas de 10 kHz.}%
\only<3->{Produto $g_1 g_2$: cada componente de $g_2$ recebe uma cópia do espectro de $g_1$.
\[
s(t) = A_c\sum_{m,n} J_m(\beta_1)\,J_n(\beta_2)\cos[(\omega_c + m\omega_1 + n\omega_2)t + m\theta_1 + n\theta_2]
\]}}

\begin{center}
\begin{tikzpicture}
\begin{axis}[espec, width=11cm, height=2.6cm, xmin=-24, xmax=24, ymax=0.85,
             xlabel={$f - f_c$ [kHz]}, ylabel={amplitude}, xtick={-20,-10,0,10,20}, ytick={0.25,0.5,0.75}]
\only<1>{\addplot+[comp, cBlue, mark options={fill=cBlue}] table[x=f, y=absJ]{dados/doistons_g1.dat};}
\only<2>{\addplot+[comp, cRed, mark options={fill=cRed}] table[x=f, y=absJ]{dados/doistons_g2.dat};}
\only<3->{\addplot+[comp, cPurple, mark options={fill=cPurple}] table[x=f, y=absJ]{dados/doistons_prod.dat};
  \node[font=\tiny, cPurple, anchor=south] at (axis cs:0,0.74) {$(0,0)$: $J_0(1)J_0(0{,}5) = 0{,}72$};
  \node[font=\tiny, cPurple, anchor=south] at (axis cs:10,0.2) {cópia em $+f_2$};
  \node[font=\tiny, cPurple, anchor=south] at (axis cs:-10,0.2) {cópia em $-f_2$};}
\end{axis}
\end{tikzpicture}
\end{center}

\onslide<4->{Surgem componentes em $f_c + m f_1 + n f_2$: intermodulação, que não aparece em AM.}

\end{frame}

% ============================================

\begin{frame}{Dois tons: componentes numéricas}

\small
Componente em $f_c + mf_1 + nf_2$ com amplitude $A_c|J_m(\beta_1)\,J_n(\beta_2)|$. $J_m(1) = 0{,}77;\ 0{,}44;\ 0{,}11$ e $J_n(0{,}5) = 0{,}94;\ 0{,}24$.

\begin{columns}[T]
\column{0.42\textwidth}
\scriptsize
\renewcommand{\arraystretch}{1.05}
\begin{tabular}{ccc}
\toprule
$(m,n)$ & $f$ [kHz] & amplitude [V] \\
\midrule
\color{cPurple}$(0,0)$ & \color{cPurple}100 & \color{cPurple}0,72 \\
\color{cBlue}$(\pm 1,0)$ & \color{cBlue}99, 101 & \color{cBlue}0,41 \\
\color{cBlue}$(\pm 2,0)$ & \color{cBlue}98, 102 & \color{cBlue}0,11 \\
\color{cRed}$(0,\pm 1)$ & \color{cRed}90, 110 & \color{cRed}0,19 \\
\color{cGreen}$(\pm 1,\pm 1)$ & \color{cGreen}89, 91, 109, 111 & \color{cGreen}0,11 \\
\color{cGreen}$(\pm 2,\pm 1)$ & \color{cGreen}88, 92, 108, 112 & \color{cGreen}0,03 \\
\bottomrule
\end{tabular}

\column{0.58\textwidth}
\begin{tikzpicture}
\begin{axis}[espec, width=6cm, height=2.9cm, xmin=86, xmax=114, ymax=0.85,
             xlabel={$f$ [kHz]}, ylabel={amplitude [V]}, xtick={90,100,110}, ytick={0.25,0.5,0.75}]
\addplot+[comp, cPurple!35, mark options={fill=cPurple!35}, restrict x to domain=86:114] table[x expr=\thisrow{f}+100, y=absJ]{dados/doistons_prod.dat};
\only<1->{\addplot+[comp, cPurple, mark options={fill=cPurple}] coordinates {(100,0.719)};}
\only<2->{\addplot+[comp, cBlue, mark options={fill=cBlue}] coordinates {(99,0.414) (101,0.414) (98,0.108) (102,0.108)};}
\only<3->{\addplot+[comp, cRed, mark options={fill=cRed}] coordinates {(90,0.185) (110,0.185)};}
\only<4->{\addplot+[comp, cGreen, mark options={fill=cGreen}] coordinates {(89,0.107) (91,0.107) (109,0.107) (111,0.107) (88,0.028) (92,0.028) (108,0.028) (112,0.028)};}
\end{axis}
\end{tikzpicture}
\end{columns}

\end{frame}

% ============================================

\begin{frame}{Muitos tons: produto de espectros}

\small
\parbox[t][1.35cm]{\linewidth}{%
\only<1>{Com $K$ tons, a fase é a soma de $K$ fases de tom único, cada uma com seu índice:
\[
\phi(t) = \sum_{k=1}^{K} \tikzmarknode{bk}{\beta_k}\sin(\omega_k t + \theta_k)
\]}%
\only<2>{Exponencial de soma vira produto de $K$ fatores; na frequência, convolução de $K$ espectros de Bessel:
\[
g(t) = \prod_{k=1}^{K} e^{j\beta_k\sin(\omega_k t + \theta_k)}
\quad\Longleftrightarrow\quad
G(f) = \tikzmarknode{conv}{G_1(f) * G_2(f) * \cdots * G_K(f)}
\]}%
\only<3->{Componentes em $f_c + \sum_k n_k f_k$, com amplitude produto de $K$ funções de Bessel:
\[
s(t) = A_c \sum_{n_1,\dots,n_K} \tikzmarknode{prod}{\Big[\prod_{k} J_{n_k}(\beta_k)\Big]} \cos\Big[\Big(\omega_c + \sum_k n_k\omega_k\Big)t + \sum_k n_k\theta_k\Big]
\]}}

\vspace{0.25cm}
\begin{center}
\begin{tikzpicture}
\begin{axis}[espec, width=11cm, height=2.4cm, xmin=-24, xmax=24, ymax=0.85,
             xlabel={$f - f_c$ [kHz]}, ylabel={amplitude}, xtick={-20,-10,0,10,20}, ytick={0.25,0.5,0.75}]
\only<1-3>{\addplot+[comp, cPurple, mark options={fill=cPurple}] table[x=f, y=absJ]{dados/doistons_prod.dat};
  \node[font=\tiny, cPurple, anchor=south] at (axis cs:0,0.74) {dois tons: portadora $0{,}72$};}
\only<4->{\addplot+[comp, cPurple!25, mark options={fill=cPurple!25}] table[x=f, y=absJ]{dados/doistons_prod.dat};
  \IfFileExists{dados/trestons_prod.dat}{\addplot+[comp, cGreen, mark options={fill=cGreen}] table[x=f, y=absJ]{dados/trestons_prod.dat};}{}
  \node[font=\tiny, cGreen, anchor=south] at (axis cs:0,0.74) {três tons: portadora $J_0(1)J_0(0{,}5)J_0(0{,}6) = 0{,}65$};}
\end{axis}
\end{tikzpicture}
\end{center}

\onslide<4->{Terceiro tom em $f_3 = 3{,}5$\,kHz com $\beta_3 = 0{,}6$: cada componente anterior ganha sua cópia, o espectro fica mais denso e a portadora perde potência.}

\only<1>{\anotacoes{
  \node[nota=cBlue, anchor=west] (n1) at ([xshift=1.2cm]bk.east) {$\beta_k = k_f A_k / f_k$: um índice por tom};
  \draw[setinha, cBlue] (n1.west) to[bend right=20] (bk.north east);
}}
\only<2>{\anotacoes{
  \node[nota=cRed, anchor=north] (n2) at ([yshift=-0.3cm]conv.south) {cada tom espalha cópias de tudo o que já existia};
  \draw[setinha, cRed] (n2.north) -- (conv.south);
}}
\only<3->{\anotacoes{
  \node[nota=cGreen, anchor=north] (n3) at ([yshift=-0.25cm]prod.south) {$|J_0(\beta_k)| < 1$: portadora cai a cada tom};
  \draw[setinha, cGreen] (n3.north) -- (prod.south);
}}

\end{frame}

% ============================================

\begin{frame}{Largura de banda}

\small
\begin{columns}[T]
\column{0.42\textwidth}
Em teoria, infinitas. Na prática, $|J_n(\beta)|$ cai rápido para $|n| > \beta + 1$.

\vspace{0.2cm}
\onslide<2->{Contando $\beta + 1$ componentes de cada lado, espaçadas de $f_m$:}
\onslide<3->{\[
\boxed{\,B_T \approx 2(\beta + 1) f_m = 2(\Delta f + f_m)\,}
\]
Regra de Carson.}

\vspace{0.2cm}
\onslide<4->{Para $\beta = 5$, a faixa guarda 99{,}4\,\% da potência.}

\column{0.58\textwidth}
\begin{tikzpicture}
\begin{axis}[espec, width=6.2cm, height=3.2cm, xmin=-9.5, xmax=9.5, ymax=0.48,
             xlabel={$n$}, ylabel={$|J_n(5)|$}, xtick={-8,-6,...,8}, ytick={0.2,0.4}]
\only<2->{\fill[cGreen, opacity=0.12] (axis cs:-6.5,0) rectangle (axis cs:6.5,0.48);}
\only<3->{\draw[cGreen, <->, thick] (axis cs:-6,0.45) -- (axis cs:6,0.45) node[midway, above, font=\scriptsize] {$B_T = 2(\beta + 1)f_m$};}
\addplot+[comp, cBlue, mark options={fill=cBlue}] table[x=n, y=absJ]{dados/comp_b5.dat};
\end{axis}
\end{tikzpicture}
\end{columns}

\end{frame}

% ============================================

\begin{frame}{Carson contra o critério de 1\,\%}

\small
\begin{columns}[T]
\column{0.58\textwidth}
\begin{tikzpicture}
\begin{axis}[espec, width=6.4cm, height=3.4cm, xmin=0, xmax=10, ymax=26,
             xlabel={$\beta$}, ylabel={$B_T/f_m$}, xtick={0,2,...,10}, ytick={0,5,10,15,20,25}]
\addplot[very thick, cBlue, domain=0:10] {2*(x+1)} node[pos=0.75, below right, font=\scriptsize] {Carson};
\onslide<2->{\addplot[thick, cRed, const plot] table[x=beta, y=crit]{dados/carson.dat};
  \node[font=\scriptsize, cRed, anchor=north west] at (axis cs:5.2,13.5) {$|J_n| \ge 0{,}01$};}
\end{axis}
\end{tikzpicture}

\column{0.42\textwidth}
Carson: reta $2(\beta + 1)$.

\vspace{0.15cm}
\onslide<2->{Contando todas as componentes acima de 1\,\% da portadora sem modulação, a banda cresce em degraus e fica acima de Carson.}

\vspace{0.15cm}
\onslide<3->{Carson guarda cerca de 98\,\% da potência: boa estimativa para dimensionar filtros e canais.}
\end{columns}

\end{frame}

% ============================================

\begin{frame}{Carson numa simulação}

\small
\parbox[t][0.95cm]{\linewidth}{%
\only<1>{FM simulado e espectro por FFT: $\Delta f = 5$\,kHz e tom de $f_m = 1$\,kHz, $\beta = 5$.}%
\only<2>{Banda de 98\,\%: abrimos a faixa a partir da portadora até juntar 98\,\% da potência. Em $\pm 5$\,kHz há só 95,9\,\%; em $\pm 6$\,kHz, 99,4\,\%.}%
\only<3->{Carson: $2(\Delta f + f_m) = 12$\,kHz, a mesma faixa que guarda 98\,\% da potência.}}

\begin{center}
\begin{tikzpicture}
\begin{axis}[espec, width=11.5cm, height=4cm, xmin=-15, xmax=15, ymin=-70, ymax=32,
             xlabel={$f - f_c$ [kHz]}, ylabel={PSD [dB]}, xtick={-15,-12,...,15}, ytick={-60,-40,-20,0}]
\only<3->{\fill[cGreen, opacity=0.12] (axis cs:-6,-70) rectangle (axis cs:6,32);
  \draw[cGreen, <->, thick] (axis cs:-6,24) -- (axis cs:6,24) node[midway, above, font=\scriptsize] {Carson, 12 kHz};}
\addplot[cBlue, thick] table[x=f, y=P]{dados/psd_tom.dat};
\only<2->{\draw[cRed, dashed] (axis cs:-5,-70) -- (axis cs:-5,4) (axis cs:5,-70) -- (axis cs:5,4)
    (axis cs:-6,-70) -- (axis cs:-6,14) (axis cs:6,-70) -- (axis cs:6,14);
  \draw[cRed!60, <->] (axis cs:-5,4) -- (axis cs:5,4) node[midway, above, font=\scriptsize] {$\pm 5$ kHz: 95,9\,\%};
  \draw[cRed, <->, thick] (axis cs:-6,14) -- (axis cs:6,14) node[midway, above, font=\scriptsize] {$\pm 6$ kHz: 99,4\,\%, passa de 98\,\%};}
\end{axis}
\end{tikzpicture}
\end{center}

\end{frame}

% ============================================

\begin{frame}{NBFM e WBFM}

\small
\begin{columns}[T]
\column{0.5\textwidth}
\centering
NBFM, $\beta = 0{,}3$

\begin{tikzpicture}
\begin{axis}[espec, width=4.8cm, height=2.4cm, xmin=-12, xmax=12, ymax=1.05,
             xlabel={$n$}, ylabel={}, xtick={-10,0,10}, ytick={0.5,1}]
\only<2->{\fill[cGreen, opacity=0.12] (axis cs:-1.5,0) rectangle (axis cs:1.5,1.05);}
\addplot+[comp, cBlue, mark options={fill=cBlue}] table[x=n, y=absJ]{dados/comp_b0p3.dat};
\end{axis}
\end{tikzpicture}

\raggedright
\onslide<2->{$B_T \approx 2f_m$, como AM: uma portadora e um par de laterais.}

\column{0.5\textwidth}
\centering
WBFM, $\beta = 10$

\begin{tikzpicture}
\begin{axis}[espec, width=4.8cm, height=2.4cm, xmin=-14, xmax=14, ymax=0.35,
             xlabel={$n$}, ylabel={}, xtick={-10,0,10}, ytick={0.1,0.3}]
\only<3->{\fill[cGreen, opacity=0.12] (axis cs:-11.5,0) rectangle (axis cs:11.5,0.35);}
\addplot+[comp, cRed, mark options={fill=cRed}] table[x=n, y=absJ]{dados/comp_b10.dat};
\end{axis}
\end{tikzpicture}

\raggedright
\onslide<3->{$B_T \approx 2\Delta f$: a banda é ditada pelo desvio, não pela mensagem.}
\end{columns}

\vspace{0.2cm}
\onslide<4->{AM ocupa $2W$ qualquer que seja a potência. FM troca banda por desempenho em ruído, que veremos no capítulo 5.}

\end{frame}

% ============================================

\begin{frame}{Banda estreita: o sinal}

\small\compacto
O sinal FM, com toda a mensagem dentro da fase:
\[
s(t) = A_c\cos\big[\tikzmarknode{wc}{\omega_c t} + \tikzmarknode{ph}{\phi(t)}\big]
\]

\vspace{0.7cm}
\onslide<2->{Cosseno da soma: \ $\cos(a+b) = \cos a\cos b - \sin a\sin b$}

\vspace{0.4cm}
\onslide<3->{Com $a = \omega_c t$ e $b = \phi(t)$:
\[
s(t) = A_c\cos(\omega_c t)\,\tikzmarknode{cphi0}{\cos\phi(t)} - A_c\sin(\omega_c t)\,\tikzmarknode{sphi0}{\sin\phi(t)}
\]}

\vspace{0.9cm}
\onslide<5->{Se $|\phi(t)| \ll 1$, quanto valem $\cos\phi$ e $\sin\phi$? A série de Taylor responde.}

\only<1->{\anotacoes{
  \node[nota=cBlue, anchor=north east] (b1) at ([xshift=-0.4cm, yshift=-0.3cm]wc.south) {portadora};
  \draw[setinha, cBlue] (b1.north east) -- (wc.south);
  \node[nota=cRed, anchor=north west] (b2) at ([xshift=0.4cm, yshift=-0.3cm]ph.south) {mensagem};
  \draw[setinha, cRed] (b2.north west) -- (ph.south);
}}
\only<4->{\anotacoes{
  \node[nota=cOrange, anchor=south] (a1) at ([xshift=-0.3cm, yshift=0.35cm]sphi0.north) {só estes dependem da mensagem};
  \draw[setinha, cOrange] (a1.south west) to[bend right=15] (cphi0.north);
  \draw[setinha, cOrange] (a1.south) -- (sphi0.north);
}}

\end{frame}

% ============================================

\begin{frame}{Banda estreita: série de Taylor}

\small\compacto
Perto de $\phi = 0$, toda função suave vira um polinômio:
\[
f(\phi) = \tikzmarknode{coef}{f(0)} + f'(0)\,\phi + \frac{f''(0)}{2!}\,\phi^2 + \frac{f'''(0)}{3!}\,\phi^3 + \cdots
\]

\vspace{0.5cm}
\begin{columns}[T]
\column{0.5\textwidth}
\renewcommand{\arraystretch}{1.2}
\centering
\begin{tabular}{c|cc|cc}
 & \onslide<2->{$\cos\phi$} & \onslide<2->{em $0$} & \onslide<3->{$\sin\phi$} & \onslide<3->{em $0$} \\
\hline
$f$    & \onslide<2->{$\cos\phi$}  & \onslide<2->{\color{cBlue}$1$}  & \onslide<3->{$\sin\phi$}  & \onslide<3->{$0$} \\
$f'$   & \onslide<2->{$-\sin\phi$} & \onslide<2->{$0$}               & \onslide<3->{$\cos\phi$}  & \onslide<3->{\color{cRed}$1$} \\
$f''$  & \onslide<2->{$-\cos\phi$} & \onslide<2->{\color{cBlue}$-1$} & \onslide<3->{$-\sin\phi$} & \onslide<3->{$0$} \\
$f'''$ & \onslide<2->{$\sin\phi$}  & \onslide<2->{$0$}               & \onslide<3->{$-\cos\phi$} & \onslide<3->{\color{cRed}$-1$} \\
\end{tabular}

\column{0.5\textwidth}
\onslide<2->{\[
\cos\phi = \tikzmarknode{c0}{\color{cBlue}1} - \frac{\phi^2}{2} + \frac{\phi^4}{24} - \cdots
\]}
\vspace{0.2cm}
\onslide<3->{\[
\sin\phi = \tikzmarknode{s1}{\color{cRed}\phi} - \frac{\phi^3}{6} + \frac{\phi^5}{120} - \cdots
\]}
\end{columns}

\vspace{0.9cm}
\onslide<4->{Derivadas em zero se alternam entre $0$ e $\pm 1$: o cosseno só tem potências pares, o seno só ímpares.}

\only<1->{\anotacoes{
  \node[nota=cPurple, anchor=north east] (t1) at ([xshift=-0.3cm, yshift=-0.3cm]coef.south) {coeficientes: derivadas em $\phi = 0$};
  \draw[setinha, cPurple] (t1.north east) -- (coef.south);
}}

\end{frame}

% ============================================

\begin{frame}{Banda estreita: truncando a série}

\small\compacto
\begin{columns}[T]
\column{0.55\textwidth}
Se $|\phi| \ll 1$, cada potência é bem menor que a anterior, e os termos em vermelho podem ser desprezados:
\[
\cos\phi = 1 - \tikzmarknode{e2}{\color{cRed}\frac{\phi^2}{2}} + \cdots \approx 1
\]
\[
\sin\phi = \phi - \tikzmarknode{e3}{\color{cRed}\frac{\phi^3}{6}} + \cdots \approx \phi
\]

\vspace{0.4cm}
\onslide<2->{Exemplo, $\phi = 0{,}2$\,rad:

\vspace{0.1cm}
\renewcommand{\arraystretch}{1.15}
\begin{tabular}{lccc}
 & exato & aproximação & erro \\
\hline
$\cos\phi$ & $0{,}980$ & $1$   & $0{,}020$ \\
$\sin\phi$ & $0{,}199$ & $0{,}2$ & $0{,}001$ \\
\end{tabular}}

\vspace{0.3cm}
\onslide<4->{Na banda estreita, $\cos\phi \approx 1$ e $\sin\phi \approx \phi$.}

\column{0.45\textwidth}
\onslide<3->{
\begin{tikzpicture}
\begin{axis}[espec, axis lines=middle, width=5cm, height=3.8cm, domain=-1.5:1.5, samples=100,
             xmin=-1.5, xmax=1.5, ymin=-1.25, ymax=1.35, xlabel={$\phi$}, ylabel={},
             xtick={-1,1}, ytick={-1,1}]
\fill[cGreen, opacity=0.15] (axis cs:-0.2,-1.25) rectangle (axis cs:0.2,1.35);
\addplot[very thick, cBlue] {cos(deg(x))};
\addplot[cBlue, dashed, thick] {1};
\addplot[very thick, cRed] {sin(deg(x))};
\addplot[cRed, dashed, thick] {x};
\node[font=\scriptsize, cBlue, anchor=south] at (axis cs:-1.05,0.55) {$\cos\phi$};
\node[font=\scriptsize, cRed, anchor=north west] at (axis cs:0.35,-0.1) {$\sin\phi$};
\node[font=\scriptsize, cGreen!70!black, anchor=north] at (axis cs:0,-0.55) {$|\phi| \le 0{,}2$};
\end{axis}
\end{tikzpicture}

\scriptsize Cheio: função exata. Tracejado: só o primeiro termo. Na faixa verde, praticamente coincidem.}
\end{columns}


\end{frame}

% ============================================

\begin{frame}{Banda estreita: a aproximação}

\small\compacto
\begin{columns}[T]
\column{0.6\textwidth}
De volta ao sinal expandido:
\[
s(t) = A_c\cos(\omega_c t)\,\tikzmarknode{cphi}{\cos\phi(t)} - A_c\sin(\omega_c t)\,\tikzmarknode{sphi}{\sin\phi(t)}
\]

\vspace{0.9cm}
\onslide<3->{Substituindo:
\[
s_{\NBFM}(t) \approx A_c\cos(\omega_c t) - A_c\,\tikzmarknode{nb}{\phi(t)\sin(\omega_c t)}
\]}

\vspace{0.2cm}
\onslide<4->{AM convencional, para comparar:
\[
s_{\AM}(t) = A_c\cos(\omega_c t) + A_c\,\tikzmarknode{am}{\mu\, m(t)\cos(\omega_c t)}
\]}

\column{0.4\textwidth}
\onslide<5->{
\centering
\begin{tikzpicture}[scale=1.15, >={Stealth}]
\draw[->, gray] (-0.2,0) -- (2.9,0);
\draw[->, very thick, cBlue] (0,0) -- (2,0) node[midway, below, font=\scriptsize] {portadora $A_c$};
\draw[->, very thick, cRed] (2,0) -- (2,0.75) node[midway, right, font=\scriptsize, align=left] {NBFM:\\ perpendicular};
\draw[->, very thick, cGreen] (2,0) -- (2.75,0) node[pos=1, below, font=\scriptsize] {AM};
\draw[cRed, dashed] (0,0) -- (2,0.75);
\draw[cRed] (0.8,0) arc (0:20.6:0.8);
\node[font=\scriptsize, cRed] at (1.05,0.18) {$\phi$};
\end{tikzpicture}

\vspace{0.1cm}
{\scriptsize No referencial da portadora: na banda estreita a parcela da mensagem gira o fasor; na AM ela estica o fasor.}}
\end{columns}

\only<2->{\anotacoes{
  \node[nota=cBlue, anchor=north] (t1) at ([yshift=-0.35cm]cphi.south) {$\approx 1$};
  \draw[setinha, cBlue] (t1.north) -- (cphi.south);
  \node[nota=cRed, anchor=north] (t2) at ([yshift=-0.35cm]sphi.south) {$\approx \phi(t)$};
  \draw[setinha, cRed] (t2.north) -- (sphi.south);
}}
\only<4->{\anotacoes{
  \node[nota=cRed, anchor=west] (a3) at ([xshift=0.25cm]nb.east) {portadora em seno};
  \draw[setinha, cRed] (a3.west) -- (nb.east);
  \node[nota=cGreen, anchor=west] (a4) at ([xshift=0.25cm]am.east) {em cosseno};
  \draw[setinha, cGreen] (a4.west) -- (am.east);
}}

\end{frame}

% ============================================

\begin{frame}{Banda estreita com um tom}

\small\compacto
Com $\phi(t) = \beta\sin(\omega_m t)$, $\beta \ll 1$:
\[
s(t) \approx A_c\cos\omega_c t - A_c\beta\,\tikzmarknode{ss}{\sin(\omega_m t)\sin(\omega_c t)}
\]

\onslide<2->{Produto de senos vira soma de cossenos:
\[
s(t) \approx A_c\cos\omega_c t + \tfrac{A_c\beta}{2}\cos(\omega_c + \omega_m)t \tikzmarknode{menos}{\color{cRed}{} - \tfrac{A_c\beta}{2}}\cos(\omega_c - \omega_m)t
\]}

\onslide<3->{\[
s_{\AM}(t) \overset{m = \cos\omega_m t}{=} A_c\cos\omega_c t + \tfrac{A_c\mu}{2}\cos(\omega_c + \omega_m)t \tikzmarknode{mais}{\color{cGreen}{} + \tfrac{A_c\mu}{2}}\cos(\omega_c - \omega_m)t
\]}

\centering
\begin{tikzpicture}
\pgfplotsset{mini/.style={espec, axis lines=middle, width=4.2cm, height=0.75cm, xmin=-1.8, xmax=1.8, ymin=-0.5, ymax=1.15,
             xtick={-1,0,1}, xticklabels={$f_c{-}f_m$, $f_c$, $f_c{+}f_m$}, ytick=\empty, ylabel={},
             x tick label style={font=\tiny, yshift=-0.15cm}}}
\onslide<2->{
\begin{axis}[mini]
\node[font=\scriptsize, cRed, anchor=south] at (axis cs:0,1.2) {banda estreita};
\addplot+[comp, cRed, mark options={fill=cRed}] coordinates {(-1,-0.35) (0,1) (1,0.35)};
\node[font=\tiny, cRed, anchor=east] at (axis cs:-1.05,-0.35) {$-\beta/2$};
\end{axis}}
\onslide<3->{
\begin{axis}[mini, at={(8cm,0)}]
\node[font=\scriptsize, cGreen, anchor=south] at (axis cs:0,1.2) {AM};
\addplot+[comp, cGreen, mark options={fill=cGreen}] coordinates {(-1,0.35) (0,1) (1,0.35)};
\node[font=\tiny, cGreen, anchor=east] at (axis cs:-1.05,0.35) {$+\mu/2$};
\end{axis}}
\onslide<4->{\node[font=\scriptsize, align=center, text width=3.1cm] at (6.1cm,0.4cm) {mesma banda, $2f_m$;\\ só muda o sinal da componente inferior};}
\end{tikzpicture}

\only<2>{\anotacoes{
  \node[nota=cBlue, anchor=west] (a1) at ([xshift=0.3cm]ss.east) {$\sin A\sin B =$\\ $\tfrac12[\cos(A-B) - \cos(A+B)]$};
  \draw[setinha, cBlue] (a1.west) -- (ss.east);
}}
\only<3->{\anotacoes{
  \node[nota=cRed, anchor=south] (a2) at ([yshift=0.3cm]menos.north) {sinal trocado};
  \draw[setinha, cRed] (a2.south) -- (menos.north);
}}

\end{frame}

% ============================================

\begin{frame}{Até onde vale a banda estreita?}

\small
\parbox[t][0.55cm]{\linewidth}{%
\only<1>{$\beta = 0{,}2$: exato e aproximado quase se confundem.}%
\only<2>{$\beta = 1$: a aproximação sai do círculo e a amplitude deixa de ser constante.}}

\begin{columns}[T]
\column{0.3\textwidth}
\centering
\begin{tikzpicture}[scale=1.3]
\draw[gray!60] (0,0) circle (1);
\draw[->, gray] (-0.2,0) -- (1.5,0);
\draw[->, very thick, cBlue] (0,0) -- (1,0);
\only<1>{\draw[very thick, cGreen] (11.5:1) arc (11.5:-11.5:1);
  \draw[very thick, cRed, dashed] (1,-0.2) -- (1,0.2);}
\only<2>{\draw[very thick, cGreen] (57.3:1) arc (57.3:-57.3:1);
  \draw[very thick, cRed, dashed] (1,-1) -- (1,1);}
\node[font=\tiny, cGreen, anchor=east] at (150:1.05) {exato: arco};
\node[font=\tiny, cRed, anchor=west] at (1.05,-0.5) {aprox.: reta};
\end{tikzpicture}

{\scriptsize Ponta do fasor no referencial da portadora.}

\column{0.7\textwidth}
\begin{tikzpicture}
\begin{axis}[onda, width=8.6cm, height=1.5cm, domain=0:2, samples=600, ymin=-1.6, ymax=1.6, ylabel={$s(t)/A_c$}]
\only<1>{\addplot[cGreen, thick] {cos(deg(2*pi*8*x + 0.2*sin(deg(2*pi*x))))};
  \addplot[cRed, dashed] {cos(deg(2*pi*8*x)) - 0.2*sin(deg(2*pi*x))*sin(deg(2*pi*8*x))};}
\only<2>{\addplot[cGreen, thick] {cos(deg(2*pi*8*x + sin(deg(2*pi*x))))};
  \addplot[cRed, dashed] {cos(deg(2*pi*8*x)) - sin(deg(2*pi*x))*sin(deg(2*pi*8*x))};}
\end{axis}
\begin{axis}[onda, width=8.6cm, height=1.2cm, domain=0:2, samples=600, ymin=-0.6, ymax=0.6, ylabel={erro}, at={(0,-2.3cm)}]
\only<1>{\addplot[cPurple] {cos(deg(2*pi*8*x + 0.2*sin(deg(2*pi*x)))) - cos(deg(2*pi*8*x)) + 0.2*sin(deg(2*pi*x))*sin(deg(2*pi*8*x))};}
\only<2>{\addplot[cPurple] {cos(deg(2*pi*8*x + sin(deg(2*pi*x)))) - cos(deg(2*pi*8*x)) + sin(deg(2*pi*x))*sin(deg(2*pi*8*x))};}
\end{axis}
\end{tikzpicture}

{\scriptsize \textcolor{cGreen}{exato}, \textcolor{cRed}{aproximado}, $f_c = 8f_m$. Regra prática: banda estreita para $\beta \lesssim 0{,}2$.}
\end{columns}

\end{frame}

% ============================================

\begin{frame}{Exemplo 1: FM com tom único}

\small
\begin{columns}[T]
\column{0.42\textwidth}
$m(t) = \cos(2\pi \cdot 5000\, t + \theta_m)$, $\Delta f = 75$\,kHz, $f_c = 100$\,MHz.

\onslide<2->{\[
\beta = \frac{\Delta f}{f_m} = \frac{75\ \text{kHz}}{5\ \text{kHz}} = 15
\]}

\onslide<3->{\[
B_T \approx 2(75 + 5) = 160\ \text{kHz}
\]}

\onslide<4->{$\beta + 1 = 16$ pares de laterais, espaçadas de 5 kHz. WBFM.}

\column{0.58\textwidth}
\onslide<2->{
\begin{tikzpicture}
\begin{axis}[espec, width=6.2cm, height=3.2cm, xmin=-22, xmax=22, ymax=0.3,
             xlabel={$n$}, ylabel={$|J_n(15)|$}, xtick={-20,-16,-8,0,8,16,20}, ytick={0.1,0.2}]
\only<3->{\fill[cGreen, opacity=0.12] (axis cs:-16.5,0) rectangle (axis cs:16.5,0.3);
  \node[font=\scriptsize, cGreen!60!black, anchor=south] at (axis cs:0,0.3) {160 kHz};}
\addplot+[comp, cBlue, mark options={fill=cBlue}, mark size=0.9pt] table[x=n, y=absJ]{dados/comp_b15.dat};
\end{axis}
\end{tikzpicture}}
\end{columns}

\end{frame}

% ============================================

\begin{frame}{Exemplo 2: rádio FM comercial}

\small
$\Delta f = 75$\,kHz, áudio até $f_m = 15$\,kHz, canais de 200\,kHz entre 88 e 108\,MHz.

\vspace{0.1cm}
\begin{columns}[T]
\column{0.4\textwidth}
\onslide<2->{\[
\beta = \frac{75}{15} = 5
\]}
\onslide<3->{\[
B_T = 2(75 + 15) = 180\ \text{kHz}
\]}
\onslide<4->{Sobram 20\,kHz de guarda por canal.}

\vspace{0.1cm}
\onslide<5->{$(108 - 88)/0{,}2 = 100$ canais.}

\column{0.6\textwidth}
\begin{tikzpicture}[x=0.024cm]
\draw[->] (-20,0) -- (330,0) node[right, font=\scriptsize] {$f$};
\foreach \c/\nome in {0/{$f_{k-1}$}, 100/{$f_k$}, 200/{$f_{k+1}$}} {
  \draw[gray] (\c-50+50,-0.08) -- (\c-50+50,0.08);
}
\foreach \c in {50,150,250} {
  \draw[gray, dashed] (\c-50,0) -- (\c-50,1.6);
}
\draw[gray, dashed] (300,0) -- (300,1.6);
\only<3->{
\foreach \c/\col in {50/cBlue!40, 150/cBlue, 250/cBlue!40} {
  \fill[\col, opacity=0.35] (\c-45,0) rectangle (\c+45,1.2);
  \node[font=\tiny, anchor=north] at (\c,0) {canal};
}
\draw[<->] (105,1.35) -- (195,1.35) node[midway, above, font=\scriptsize] {180 kHz};}
\only<4->{\draw[<->, cRed] (100,0.6) -- (105,0.6);
  \node[nota=cRed, anchor=south] at (100,1.65) {guarda};
  \draw[setinha, cRed] (100,1.65) -- (102,0.65);}
\draw[<->] (100,-0.55) -- (200,-0.55) node[midway, below, font=\scriptsize] {200 kHz};
\end{tikzpicture}
\end{columns}

\end{frame}
